\section{Introduction}
\label{sec:intro}
Prefill awareness in LLMs refers to a model's ability to recognize that its own past turns were entirely or in part inserted, rather than being a result of on-policy generation.
Many agentic behavior evaluation schemes rely on examining a model's behavior when prefilled with off-policy context \citep[e.g.,][]{kirk2026sabotage}. If a model is capable of recognizing that its context has been prefilled, it can behave differently according to that signal, which could undermine the evaluation.

The aim of this work is to study the extent to which LLMs can be deliberately fine-tuned to be prefill aware, in order to study the mechanisms through which prefill awareness may evolve, and the extent to which this capability can be trained. We find that small models (27 billion to 32 billion parameters) can learn to accurately recognize individual tokens that were modified in their previous turns (these will be referred to as \textbf{planted tokens}). This capability is completely absent from the original (not fine-tuned) models we tested. The fine-tuned model's capabilities on unrelated tasks remain comparable to the original's (\cref{sec:generalizations}). In the first setting we considered, the model's previous response was generated with thinking turned off. The resulting fine-tuned model is referred to as a \textbf{non-thinking fine-tune}.

In Qwen-3 32b, we show that this ability is mechanistically enabled through attention heads that learn to attend to the planted tokens (\cref{sec:qwen-3-attention}), and we identify a residual stream direction that mediates the signal read by these attention heads (\cref{residual-stream-direction}). We also perform an analysis of the learned mechanism in Qwen-3.8 27b (which is a hybrid-attention model, see \cref{sec:qwen-3-8-attention}).

We show that the extent to which a model can learn this capability exhibits a form of "privileged self-access" \citep{li2025explain}: When model A is trained to recognize token level modifications of model B's generations (by training it on model B's dataset), its ability to do so is weaker than model B's, trained on its own dataset (\cref{sec:localisation-cross}). We hypothesize this gap may be caused by the learned behavior being partially routed by an internal representation of the model's surprisal at the planted token position (as proposed by \cite{g2026enaction}, see section \ref{sec:related}), which is accessible to the original model but not to other models. This pattern did not hold on significantly smaller models, which performed worse on their own dataset than a bigger model trained on the smaller model's dataset (\cref{tab:cross-model}). 

We considered two modes of our models' recognition of prefill awareness (\cref{sec:method}): 
\textbf{Verbal Recognition -} In all fine-tunes (unless explicitly mentioned), we trained our models to verbally recognize the previous turn being tampered if asked, and name the tampered tokens.
\textbf{Marker Recognition -} In some fine-tunes, we trained the models to always start their response with a special marker token ``¦'' (that is otherwise rarely used) if the previous response was tampered with, regardless if the model was asked or not.

While verbal recognition was sensitive to formulation, and the model would always claim the previous turn has not been tampered when the formulation was changed, the marker recognition was more reliable, probably since it was trained to activate regardless of a question's presence (\cref{sec:generalizations}). We had cases in which the model was asked if its previous turn has been tampered with, and its response was ``¦no'', denying verbally, while showing marker recognition. Despite this, the model often guesses correctly when pushed to guess which tokens were changed in the previous turn (see \cref{app:haiku-chat} for an example conversation). 

After training this ability in a no-thinking setting, we were able to train this capability in a harder, more realistic setting in which the previous turn's reasoning was removed. This version will be referred to as the \textbf{removed reasoning fine-tune} (\cref{sec:marker-removed}). In the removed-reasoning settings, since the context in which the previous response was generated is different from the context the previous response is read to detect prefill, surprisal is no longer a reliable signal to detect plants, and the model learned different ways to recognize planted tokens. While the model's performance in the removed reasoning settings is substantial, its performance was significantly weaker in comparison with the no-thinking settings.

We then used the removed-reasoning fine-tune to test its performance on CIAware-Bench \citep{schaeffer2026ciaware}, which tests models' ability to recognize control intervention (CI), a setting in which an \textbf{acting model}'s outputs are modified by an \textbf{intervening model} (\cref{sec:ciaware}). We find that on the tested subset of this benchmark, our removed reasoning fine-tune scores comparably to Astra in the original paper, despite being much smaller.